\documentclass[aspectratio=169]{beamer}
\usetheme{Madrid}
\usecolortheme{default}

\usepackage{amsmath,amssymb,amsfonts,bm,mathtools}
\usepackage{physics}
\usepackage{braket}

\title{From Second Quantization and Quasispin to Pauli Strings}
\author{Morten Hjorth-Jensen}
\date{2026}

\begin{document}

\frame{\titlepage}

\begin{frame}{Scope of these notes}
These notes explain how to rewrite many-body Hamiltonians in terms of Pauli strings
built from
\[
I,\qquad X,\qquad Y,\qquad Z,
\]
so that they are ready for quantum-computing implementations.

\medskip

Main topics:
\begin{itemize}
\item the Lipkin Hamiltonian in quasispin form,
\item detailed mappings to Pauli operators,
\item compressed and direct qubit encodings,
\item second-quantized fermionic operators \(a^\dagger,a\),
\item the nuclear pairing Hamiltonian as a second worked example,
\item Jordan--Wigner-type mappings from fermions to qubits.
\end{itemize}
\end{frame}

\begin{frame}{Notation}
We use:
\begin{itemize}
\item \(a_p^\dagger\): fermion creation operator in orbital \(p\),
\item \(a_p\): fermion annihilation operator in orbital \(p\),
\item \(n_p=a_p^\dagger a_p\): number operator,
\item \(X_i,Y_i,Z_i\): Pauli operators on qubit \(i\),
\item \(I_i\): identity on qubit \(i\).
\end{itemize}

Throughout, we avoid \(c_p^\dagger,c_p\) notation and use only
\[
a_p^\dagger,\qquad a_p.
\]

For quasispin systems we use collective operators
\[
J_x,\; J_y,\; J_z,\qquad J_\pm = J_x \pm i J_y.
\]
\end{frame}

\begin{frame}{Why Pauli decompositions matter}
Quantum computers natively implement unitary evolutions generated by sums of Pauli strings,
for example
\[
H = \sum_\alpha h_\alpha P_\alpha,
\qquad
P_\alpha \in \{I,X,Y,Z\}^{\otimes n}.
\]

Once a Hamiltonian is written in this form, one can:
\begin{itemize}
\item estimate energies in VQE,
\item generate Trotterized time evolution,
\item build QAOA circuits,
\item analyze measurement grouping and commutativity.
\end{itemize}

So the essential preprocessing step is:
\begin{quote}
rewrite the many-body Hamiltonian as a weighted sum of Pauli strings.
\end{quote}
\end{frame}

\section{Part I: Lipkin Hamiltonian in quasispin form}

\begin{frame}{Lipkin Hamiltonian}
We consider the Lipkin Hamiltonian in quasispin representation
\[
H
=
\frac{\epsilon}{2}\left(J_z-\frac{N}{2}\right)
+
\frac{V}{2}\left(J_+^2+J_-^2\right)
+
\frac{W}{2}J_z^2.
\]

Here:
\begin{itemize}
\item \(N\) is the number of particles,
\item \(\epsilon\) sets the one-body splitting,
\item \(V\) controls pair scattering between the two levels,
\item \(W\) weights the quadratic \(J_z^2\) term.
\end{itemize}
\end{frame}

\begin{frame}{Quasispin algebra}
The collective operators satisfy the \(su(2)\) algebra
\[
[J_x,J_y]=iJ_z,\qquad
[J_y,J_z]=iJ_x,\qquad
[J_z,J_x]=iJ_y,
\]
or, equivalently,
\[
[J_z,J_\pm]=\pm J_\pm,
\qquad
[J_+,J_-]=2J_z.
\]

The basis states are
\[
\ket{j,m},
\qquad
m=-j,-j+1,\dots,j,
\]
with
\[
J_z\ket{j,m}=m\ket{j,m}.
\]
\end{frame}

\begin{frame}{Action of ladder operators}
The ladder operators act as
\[
J_+\ket{j,m}
=
\sqrt{j(j+1)-m(m+1)}\ket{j,m+1},
\]
\[
J_-\ket{j,m}
=
\sqrt{j(j+1)-m(m-1)}\ket{j,m-1}.
\]

Therefore
\[
J_+^2\ket{j,m}
=
\sqrt{j(j+1)-m(m+1)}
\sqrt{j(j+1)-(m+1)(m+2)}
\ket{j,m+2},
\]
and analogously for \(J_-^2\).
\end{frame}

\begin{frame}{Useful identity for the interaction term}
Using
\[
J_\pm = J_x \pm iJ_y,
\]
we compute
\[
J_+^2 + J_-^2
=
(J_x+iJ_y)^2 + (J_x-iJ_y)^2
=
2(J_x^2-J_y^2).
\]

Hence the Lipkin Hamiltonian can also be written as
\[
H
=
\frac{\epsilon}{2}\left(J_z-\frac{N}{2}\right)
+
V(J_x^2-J_y^2)
+
\frac{W}{2}J_z^2.
\]

This form is often the most convenient when collective spin operators are mapped directly to qubits.
\end{frame}

\begin{frame}{Two mapping strategies}
There are two main strategies for turning the Lipkin Hamiltonian into Pauli strings:
\begin{enumerate}
\item \textbf{Compressed encoding of the quasispin basis}: encode the \((2j+1)\)-dimensional quasispin space into \(\lceil \log_2(2j+1)\rceil\) qubits.
\item \textbf{Direct collective-spin mapping}: write
\[
J_\alpha = \frac12\sum_i \sigma_i^\alpha
\]
on a multi-qubit register and expand the Hamiltonian.
\end{enumerate}

We will use:
\begin{itemize}
\item compressed encoding for \(N=4\),
\item direct collective mapping for \(N=6\).
\end{itemize}
\end{frame}

\subsection{Case \(N=4\): compressed 3-qubit encoding}

\begin{frame}{Four particles: quasispin dimension and encoding}
For \(N=4\),
\[
j=\frac{N}{2}=2,
\qquad
2j+1 = 5.
\]

Thus the quasispin sector has dimension \(5\), which fits into \(3\) qubits since
\[
2^2=4<5\le 8=2^3.
\]

We choose the binary map
\[
\ket{2,-2}\leftrightarrow \ket{000},\quad
\ket{2,-1}\leftrightarrow \ket{001},\quad
\ket{2,0}\leftrightarrow \ket{010},
\]
\[
\ket{2,1}\leftrightarrow \ket{011},\quad
\ket{2,2}\leftrightarrow \ket{100}.
\]

The states \(\ket{101},\ket{110},\ket{111}\) are unphysical.
\end{frame}

\begin{frame}{Projectors and transition operators on qubits}
For one qubit:
\[
\ket{0}\bra{0}=\frac{I+Z}{2},
\qquad
\ket{1}\bra{1}=\frac{I-Z}{2},
\]
\[
\ket{1}\bra{0}=\frac{X-iY}{2},
\qquad
\ket{0}\bra{1}=\frac{X+iY}{2}.
\]

Therefore:
\begin{itemize}
\item diagonal operators are expanded with \(I\) and \(Z\),
\item off-diagonal transitions are expanded with \(X\) and \(Y\),
\item any finite operator in the encoded space can be written as a sum of Pauli strings.
\end{itemize}
\end{frame}

\begin{frame}{Three-qubit form of \(J_z\)}
Since
\[
J_z = \sum_{m=-2}^2 m\ket{m}\bra{m},
\]
we obtain
\[
J_z
=
-2\ket{000}\bra{000}
-\ket{001}\bra{001}
+\ket{011}\bra{011}
+2\ket{100}\bra{100}.
\]

Expanding the projectors gives
\[
J_z
=
-\frac{1}{4} IZI
+\frac{1}{4} IZZ
-\frac{1}{2} ZII
-\frac{1}{2} ZIZ
-\frac{3}{4} ZZI
-\frac{1}{4} ZZZ.
\]
\end{frame}

\begin{frame}{Three-qubit form of \(J_z^2\)}
Similarly,
\[
J_z^2 = \sum_{m=-2}^2 m^2 \ket{m}\bra{m},
\]
which yields
\[
J_z^2
=
\frac{5}{4} III
+\frac{3}{4} IIZ
+IZI
+IZZ
+\frac{1}{4} ZII
-\frac{1}{4} ZIZ.
\]

These formulas are exact on the physical 5-dimensional subspace.
\end{frame}

\begin{frame}{Three-qubit form of the ladder terms}
For \(j=2\),
\[
J_+\ket{-2}=2\ket{-1},\qquad
J_+\ket{-1}=\sqrt{6}\ket{0},
\]
\[
J_+\ket{0}=\sqrt{6}\ket{1},\qquad
J_+\ket{1}=2\ket{2}.
\]

Therefore
\[
J_+^2+J_-^2
=
\frac{3+\sqrt{6}}{2}\,IXI
+\frac{\sqrt{6}-3}{2}\,IXZ
+\frac{\sqrt{6}}{2}\,XXI
+\frac{\sqrt{6}}{2}\,XXZ
\]
\[
\qquad
+\frac{\sqrt{6}}{2}\,YYI
+\frac{\sqrt{6}}{2}\,YYZ
+\frac{3+\sqrt{6}}{2}\,ZXI
+\frac{\sqrt{6}-3}{2}\,ZXZ.
\]
\end{frame}

\begin{frame}{Final \(N=4\), 3-qubit Lipkin Hamiltonian}
Substituting into
\[
H = \frac{\epsilon}{2}(J_z-2) + \frac{V}{2}(J_+^2+J_-^2) + \frac{W}{2}J_z^2
\]
gives
\[
H_{N=4}^{(3q)}
=
\left(-\epsilon+\frac{5W}{8}\right)III
+\frac{3W}{8}IIZ
+\left(-\frac{\epsilon}{8}+\frac{W}{2}\right)IZI
+\left(\frac{\epsilon}{8}+\frac{W}{2}\right)IZZ
\]
\[
+\left(-\frac{\epsilon}{4}+\frac{W}{8}\right)ZII
+\left(-\frac{\epsilon}{4}-\frac{W}{8}\right)ZIZ
-\frac{3\epsilon}{8}ZZI
-\frac{\epsilon}{8}ZZZ
\]
\[
+\frac{V(3+\sqrt{6})}{4}IXI
+\frac{V(\sqrt{6}-3)}{4}IXZ
+\frac{V\sqrt{6}}{4}XXI
+\frac{V\sqrt{6}}{4}XXZ
+\frac{V\sqrt{6}}{4}YYI
+\frac{V\sqrt{6}}{4}YYZ
+\frac{V(3+\sqrt{6})}{4}ZXI
+\frac{V(\sqrt{6}-3)}{4}ZXZ.
\]
\end{frame}

\begin{frame}{Code-space penalty}
Because the 3-qubit register has 8 basis states but only 5 are physical, one often adds
\[
H_{\mathrm{pen}} = \lambda P_{\mathrm{unphys}},
\]
where
\[
P_{\mathrm{unphys}}
=
\ket{101}\bra{101}
+\ket{110}\bra{110}
+\ket{111}\bra{111}.
\]

This energetically suppresses leakage outside the encoded quasispin space.
\end{frame}

\subsection{Case \(N=6\): direct 6-qubit mapping}

\begin{frame}{Six particles: direct collective-spin realization}
For \(N=6\), a simple and transparent mapping uses six qubits and defines
\[
J_x = \frac{1}{2}\sum_{i=1}^{6} X_i,
\qquad
J_y = \frac{1}{2}\sum_{i=1}^{6} Y_i,
\qquad
J_z = \frac{1}{2}\sum_{i=1}^{6} Z_i.
\]

This is not qubit-optimal, since the symmetric quasispin sector has dimension \(7\), but it leads to a very clean Pauli-string Hamiltonian.
\end{frame}

\begin{frame}{Expanding the collective squares}
Using \(X_i^2=Y_i^2=Z_i^2=I\),
\[
J_x^2
=
\frac{1}{4}\sum_{i=1}^{6} I
+
\frac{1}{4}\sum_{i\neq j} X_iX_j
=
\frac{6}{4}I + \frac{1}{2}\sum_{i<j} X_iX_j,
\]
and similarly
\[
J_y^2
=
\frac{6}{4}I + \frac{1}{2}\sum_{i<j} Y_iY_j,
\]
\[
J_z^2
=
\frac{6}{4}I + \frac{1}{2}\sum_{i<j} Z_iZ_j.
\]
\end{frame}

\begin{frame}{Interaction term on 6 qubits}
Using
\[
J_+^2+J_-^2=2(J_x^2-J_y^2),
\]
we get
\[
\frac{V}{2}(J_+^2+J_-^2)
=
V(J_x^2-J_y^2)
=
\frac{V}{2}\sum_{i<j}(X_iX_j-Y_iY_j).
\]

Thus the interaction is an all-to-all two-body Pauli interaction.
\end{frame}

\begin{frame}{Linear and quadratic \(J_z\) terms on 6 qubits}
For the linear term,
\[
\frac{\epsilon}{2}\left(J_z-\frac{N}{2}\right)
=
\frac{\epsilon}{2}\left(\frac12\sum_{i=1}^{6}Z_i - 3I\right)
=
\frac{\epsilon}{4}\sum_{i=1}^{6} Z_i - \frac{3\epsilon}{2}I.
\]

For the quadratic term,
\[
\frac{W}{2}J_z^2
=
\frac{W}{2}\left(\frac{6}{4}I+\frac12\sum_{i<j}Z_iZ_j\right)
=
\frac{3W}{4}I+\frac{W}{4}\sum_{i<j}Z_iZ_j.
\]
\end{frame}

\begin{frame}{Final \(N=6\), 6-qubit Lipkin Hamiltonian}
Collecting all terms gives
\[
H_{N=6}^{(6q)}
=
\left(-\frac{3\epsilon}{2}+\frac{3W}{4}\right)I
+
\frac{\epsilon}{4}\sum_{i=1}^{6} Z_i
+
\frac{V}{2}\sum_{i<j}(X_iX_j-Y_iY_j)
+
\frac{W}{4}\sum_{i<j} Z_iZ_j.
\]

This is already directly usable in qubit-based quantum algorithms.
\end{frame}

\begin{frame}{Comparison of the two Lipkin mappings}
\begin{center}
\begin{tabular}{lcc}
\hline
Case & Qubits & Main feature \\
\hline
\(N=4\), compressed & 3 & minimal encoding \\
\(N=6\), direct & 6 & simplest Pauli structure \\
\hline
\end{tabular}
\end{center}

Compressed mappings save qubits but require code-space control.
Direct mappings use more qubits but give very transparent Hamiltonians.
\end{frame}

\section{Part II: Fermions, second quantization, and Pauli matrices}

\begin{frame}{Second quantization for fermions}
In fermionic many-body theory, the basic operators are
\[
a_p^\dagger,\qquad a_p,
\]
satisfying anticommutation relations
\[
\{a_p,a_q^\dagger\}=\delta_{pq},
\qquad
\{a_p,a_q\}=0,
\qquad
\{a_p^\dagger,a_q^\dagger\}=0.
\]

A general second-quantized Hamiltonian has the form
\[
H
=
\sum_{pq} h_{pq} a_p^\dagger a_q
+
\frac{1}{4}\sum_{pqrs} v_{pqrs} a_p^\dagger a_q^\dagger a_s a_r.
\]
\end{frame}

\begin{frame}{Why fermion-to-qubit mappings are needed}
A quantum computer acts on qubits, not directly on fermionic Fock states.
So we need a mapping
\[
a_p^\dagger,\; a_p,\; n_p=a_p^\dagger a_p
\quad \longrightarrow \quad
\text{Pauli strings}.
\]

The most common constructions are:
\begin{itemize}
\item Jordan--Wigner,
\item Bravyi--Kitaev,
\item parity mapping.
\end{itemize}

We focus here on Jordan--Wigner, since it is the most direct mathematically.
\end{frame}

\begin{frame}{Jordan--Wigner transformation}
For orbital \(p\), the Jordan--Wigner mapping is
\[
a_p^\dagger
=
\left(\prod_{j=0}^{p-1} Z_j\right)\frac{X_p-iY_p}{2},
\qquad
a_p
=
\left(\prod_{j=0}^{p-1} Z_j\right)\frac{X_p+iY_p}{2}.
\]

The string of \(Z\)'s ensures the correct fermionic anticommutation relations.
\end{frame}

\begin{frame}{Number operators under Jordan--Wigner}
Using the mapping above,
\[
n_p = a_p^\dagger a_p = \frac{1-Z_p}{2}.
\]

This is one of the simplest and most useful identities:
\[
n_p \leftrightarrow \frac{1-Z_p}{2}.
\]

Thus one-body diagonal terms in fermionic Hamiltonians immediately become sums of \(I\) and \(Z\).
\end{frame}

\begin{frame}{Hopping terms under Jordan--Wigner}
For \(p<q\), one finds that
\[
a_p^\dagger a_q + a_q^\dagger a_p
\]
maps to strings of the form
\[
\frac12
\left(
X_p Z_{p+1}\cdots Z_{q-1} X_q
+
Y_p Z_{p+1}\cdots Z_{q-1} Y_q
\right).
\]

So hopping terms become nonlocal Pauli strings because of the fermionic parity chain.
\end{frame}

\section{Part III: Nuclear pairing Hamiltonian}

\begin{frame}{Nuclear pairing Hamiltonian}
A standard reduced pairing Hamiltonian is
\[
H_{\mathrm{pair}}
=
\sum_{p} \epsilon_p \, N_p
-
G \sum_{p,q} P_p^\dagger P_q,
\]
where
\[
N_p = a_p^\dagger a_p + a_{\bar p}^\dagger a_{\bar p},
\]
\[
P_p^\dagger = a_p^\dagger a_{\bar p}^\dagger,
\qquad
P_p = a_{\bar p} a_p.
\]

Here \(p\) and \(\bar p\) denote time-reversed partners.
\end{frame}

\begin{frame}{Pairing operators and quasispin structure}
The pairing Hamiltonian is naturally associated with a quasispin algebra:
\[
S_p^+ = P_p^\dagger,
\qquad
S_p^- = P_p,
\qquad
S_p^z = \frac12(N_p-1).
\]

These satisfy
\[
[S_p^+,S_q^-]=2\delta_{pq} S_p^z,
\qquad
[S_p^z,S_q^\pm]=\pm \delta_{pq}S_p^\pm.
\]

Thus pairing and Lipkin Hamiltonians are closely related at the algebraic level:
both can be expressed in terms of collective \(su(2)\)-type generators.
\end{frame}

\begin{frame}{One pair level: direct qubit mapping}
For a single pair \((p,\bar p)\), the two-dimensional paired subspace is naturally mapped to one qubit:
\[
\ket{0}_p \equiv \text{pair unoccupied},
\qquad
\ket{1}_p \equiv \text{pair occupied}.
\]

In this reduced pairing subspace, the quasispin operators become
\[
S_p^+ = \frac{X_p-iY_p}{2},
\qquad
S_p^- = \frac{X_p+iY_p}{2},
\qquad
S_p^z = -\frac{Z_p}{2}.
\]

Also,
\[
N_p = 1 - Z_p.
\]
\end{frame}

\begin{frame}{Pairing Hamiltonian in pair-qubit form}
Substituting
\[
N_p = 1-Z_p,
\qquad
P_p^\dagger = \frac{X_p-iY_p}{2},
\qquad
P_p = \frac{X_p+iY_p}{2},
\]
we obtain
\[
H_{\mathrm{pair}}
=
\sum_p \epsilon_p (1-Z_p)
-
G\sum_{p,q} S_p^+ S_q^-.
\]

Since
\[
S_p^+S_q^- + S_q^+S_p^-
=
\frac12(X_pX_q + Y_pY_q),
\]
the interaction becomes an \(XY\)-type coupling between pair qubits.
\end{frame}

\begin{frame}{Explicit pairing Hamiltonian for \(\Omega\) pair levels}
For \(\Omega\) pair orbitals, the pairing Hamiltonian in the paired subspace can be written as
\[
H_{\mathrm{pair}}
=
\sum_{p=1}^{\Omega} \epsilon_p (1-Z_p)
-
\frac{G}{2}\sum_{p\ne q}(X_pX_q+Y_pY_q)
-
G\sum_p S_p^+S_p^-.
\]

Since
\[
S_p^+S_p^- = \frac{1-Z_p}{2},
\]
this yields
\[
H_{\mathrm{pair}}
=
\sum_{p=1}^{\Omega}\left(\epsilon_p-\frac{G}{2}\right)(1-Z_p)
-
\frac{G}{2}\sum_{p\ne q}(X_pX_q+Y_pY_q).
\]

This is already in Pauli-string form.
\end{frame}

\begin{frame}{Example: two pair levels}
For two pair orbitals \(p=1,2\), the paired-subspace Hamiltonian becomes
\[
H_{\mathrm{pair}}^{(2)}
=
\left(\epsilon_1-\frac{G}{2}\right)(1-Z_1)
+
\left(\epsilon_2-\frac{G}{2}\right)(1-Z_2)
-
\frac{G}{2}(X_1X_2+Y_1Y_2).
\]

This is one of the simplest nontrivial pairing Hamiltonians ready for quantum simulation.
\end{frame}

\begin{frame}{What if one uses full fermionic orbitals instead?}
If one keeps both \(p\) and \(\bar p\) as separate fermionic modes and applies Jordan--Wigner directly, then
\[
a_p^\dagger a_{\bar p}^\dagger,\qquad
a_{\bar p}a_p
\]
become longer Pauli strings with parity \(Z\)-chains.

This is fully general and exact in the full Fock space, but less compact.
The pair-qubit reduction is much more efficient when seniority-zero paired states are the physically relevant sector.
\end{frame}

\begin{frame}{Lipkin versus pairing: structural comparison}
Both Hamiltonians admit spin-like descriptions:
\begin{itemize}
\item Lipkin:
\[
H_{\mathrm{Lipkin}}
=
\frac{\epsilon}{2}\left(J_z-\frac{N}{2}\right)
+\frac{V}{2}(J_+^2+J_-^2)
+\frac{W}{2}J_z^2,
\]
\item Pairing:
\[
H_{\mathrm{pair}}
=
\sum_p \epsilon_p N_p - G\sum_{p,q} P_p^\dagger P_q.
\]
\end{itemize}

Both become sums of Pauli strings after:
\begin{itemize}
\item choosing an encoding,
\item expressing collective or pair operators in terms of \(X,Y,Z\).
\end{itemize}
\end{frame}

\begin{frame}{General recipe: from second quantization to qubits}
For a fermionic Hamiltonian:
\begin{enumerate}
\item write the Hamiltonian in terms of \(a_p^\dagger,a_p\),
\item decide whether a reduced quasispin or pair-qubit subspace is appropriate,
\item otherwise apply Jordan--Wigner or a related mapping,
\item simplify number and pair operators,
\item collect the resulting Pauli strings.
\end{enumerate}

For collective-spin models:
\begin{enumerate}
\item write the model in \(J_x,J_y,J_z\) or ladder-operator form,
\item choose compressed or direct qubit encoding,
\item expand diagonal and off-diagonal matrix elements into Pauli strings.
\end{enumerate}
\end{frame}

\begin{frame}{Summary}
We updated the Lipkin notes by:
\begin{itemize}
\item adding more mathematical details,
\item using \(a^\dagger,a\) notation consistently,
\item explaining both compressed and direct qubit mappings,
\item adding the nuclear pairing Hamiltonian as a second detailed example,
\item showing how second-quantized Hamiltonians become Pauli-string Hamiltonians ready for quantum computation.
\end{itemize}

The main conceptual bridge is:
\[
\text{many-body algebra}
\quad \longrightarrow \quad
\text{spin/qubit algebra}
\quad \longrightarrow \quad
\text{Pauli strings}.
\]
\end{frame}

\begin{frame}{End}
\centering
{\Large Questions?}
\end{frame}

\end{document}
